\subsection{Algebraically free families. Pure extensions}

Let us recall (IV, p.~4) the following definition :

DEFINITION 1. --- Let E be an extension of a field K; a family \(\mathbf{x} = (x_i)_{i \in I}\) of elements of E is said to be algebraically free over K if the monomials \(\mathbf{x}^\alpha = \prod_{i \in I} \mathbf{x}_i^{\alpha_i}\) with respect to the \(x_i\) (for \(\alpha = (\alpha_i)_{i \in I}\) in \(\mathbf{N}^{(I)}\) ) are linearly independent over K. In the contrary case the family is said to be algebraically related over K. Definition 1 may also be expressed as follows:

PROPOSITION 1. --- For a family \((x_{i})_{i\in I}\) of elements of an extension E of a field K to be algebraically free over K it is necessary and sufficient for the relation \(f((x_i)) = 0\) , where \(f\) is a polynomial in \(\mathbf{K}[\mathbf{X}_i]_{i\in I}\) to imply \(f = 0\) .

DEFINITION 2. --- Let E be an extension of a field K. A family \((x_{i})_{i\in I}\) of elements of E is called a pure basis of E (over K) if it is algebraically free and \(\mathrm{E}=\mathrm{K}(x_{i})_{i\in I}\) . The extension E of K is also called pure if it possesses a pure basis.

The empty family is algebraically free, hence K is a pure extension of itself. With the notation of Def. 2, every element \(x_{i}\) is transcendental over K; if I is non-empty, E is thus a transcendental extension of K.

PROPOSITION 2. --- Let E and E' be two fields and u an isomorphism of a subfield K of E onto a subfield K' of E'. Let \(\mathbf{x} = (x_i)_{i \in I}\) (resp. \(\mathbf{x}' = (x_i')_{i \in I}\) ) be a family of elements of E (resp. E') which is algebraically free over K (resp. K'). Then there exists a unique isomorphism v of L = K(x\_i)\_\{i ∈ I\} onto L' = K'(x\_i')\_\{i ∈ I\} which induces u on K and maps x\_i to x\_i' for each i ∈ I.

The uniqueness of \(v\) is clear. Let us put \(\mathbf{A} = \mathbf{K}[x_i]_{i \in I}\) and \(\mathbf{A}' = \mathbf{K}'[x_i']_{i \in I}\) . By hypothesis the monomials \(\mathbf{x}^\alpha = \prod_{i \in I} x_i^{\alpha_i}\) (for \(\alpha = (\alpha_i)_{i \in I}\) in \(\mathbf{N}^{(I)}\) ) form a basis of the

vector K-space A and there is a similar property of \(\mathbf{A}'\) . Hence there exists a ring isomorphism \(w: \mathbf{A} \to \mathbf{A}'\) mapping each element \(\sum_{\alpha \in \mathbf{N}^{(I)}} c_{\alpha} x^{\alpha}\) to \(\sum_{\alpha \in \mathbf{N}^{(I)}} u(c_{\alpha}) x'^{\alpha}\) . Since

L is the field of fractions of A and L' that of A', the isomorphism w extends to an isomorphism v of the field L onto the field L'.

COROLLARY. --- For an extension E of a field K to be pure it is necessary and sufficient for E to be K-isomorphic to a field of rational fractions over K. More precisely, if the family \((x_{i})_{i\in I}\) is a pure basis of E, then there exists a unique K-isomorphism of \(\mathrm{K}(\mathbf{X}_{i})_{i\in I}\) onto E which maps \(X_{i}\) to \(x_{i}\) for each \(i\in I\) .

Remark. --- It is clear that in an extension E of K an algebraically free family over K consists of linearly independent elements over K (hence pairwise distinct); in other words, it is also a free family for the vector space structure of E (with respect to K). But the converse is false, for if E is an algebraic extension of K, any non-empty family of elements of E (and a fortiori a non-empty family of linearly independent elements over K) is never algebraically free over K. When there is a risk of confusion, we shall say that a subset of an extension E of K which is free for the vector space structure of E with respect to K is linearly free over K.

Let E be an extension of a field K. A subset S of E is said to be algebraically free (over K) if the family defined by the identity mapping of S onto itself is algebraically free. The elements of an algebraically free subset of E are also called algebraically independent. If a subset of E is not algebraically free, it is said to be algebraically related and that its elements are algebraically dependent. For a family \((x_{i})_{i\in I}\) of elements of E to be algebraically free it is necessary and sufficient that \(i\mapsto x_{i}\) should be a bijection of I onto an algebraically free subset of E.

Every subset of an algebraically free set is algebraically free. Furthermore:

PROPOSITION 3. --- For a family \((x_{i})_{i\in I}\) of elements of an extension E of a field K to be algebraically free over K it is necessary and sufficient that every finite subfamily of \((x_{i})_{i\in I}\) should be algebraically free over K.

This proposition follows immediately from Def. 1.

